\section{Allgemeines}\label{sec:allgemeines}

Dieser Abschnitt enthält die grundlegenden Angaben zum Tarif, zu den
Beobachtungseinheiten, zur Altersbestimmung und zu den einbezogenen
Rahmen-Berechnungsgrundlagen.

\subsection{Allgemeine Darstellung des Tarifs}\label{sec:tarifdarstellung}

\textbf{Tarifname:} VitalPlus (unverändert seit Einführung zum
01.01.1998; kein abweichender Verkaufsname).

\textbf{Versicherungsform:} Krankheitskostenvollversicherung für die
Leistungsbereiche ambulante Heilbehandlung (mit absolutem Selbstbehalt
von 0 bis 100~EUR p.\,a.), stationäre Heilbehandlung (allgemeine
Krankenhausleistungen, Mehrbettzimmer) sowie Zahnbehandlung und
Zahnersatz (Erstattungssatz Zahnersatz 70--80\,\%). Einzelversicherung;
eine Gruppenversicherung wird im Tarif VitalPlus nicht geführt.

\textbf{Charakterisierung im Sinne des DAV-Hinweises zu TB:}
\begin{itemize}\itemsep2pt
  \item Der Tarif ist \textbf{substitutiv} im Sinne des \S\,146 VAG und
        nach \textbf{Art der Lebensversicherung} kalkuliert.
  \item Der Tarif ist \textbf{arbeitgeberzuschussfähig} gemäß
        \S\,257 SGB\,V.
  \item Der \textbf{gesetzliche Beitragszuschlag} nach \S\,149 VAG wird
        erhoben (Abschnitt~\ref{sec:gz}).
  \item Der Tarif ist \textbf{geschlechtsabhängig (Bisex)} kalkuliert und
        gehört zur \textbf{alten Welt}: Er wurde mit Einführung des
        Übertragungswertes zum \textbf{31.12.2008 für das Neugeschäft
        geschlossen}; alle Verträge wurden vor dem 01.01.2009
        abgeschlossen, ein Anspruch auf Mitgabe eines
        Übertragungswertes nach \S\,146 Abs.\,1 Nr.\,5 VAG besteht
        nicht (vgl. Abschnitt~\ref{sec:uew}). Die vorliegende TB dient
        der Bestandsführung, der Beitragsanpassung und individuellen
        Tarifänderungen.
  \item Es handelt sich \textbf{nicht} um einen Bausteintarif.
\end{itemize}

\textbf{Versicherbarer Personenkreis:} In Deutschland lebende Personen,
die nicht in der gesetzlichen Krankenversicherung pflichtversichert
sind; Aufnahme von Kindern und Jugendlichen im Rahmen der
Kindernachversicherung ohne Risikoprüfung.

\textbf{Allgemeine Versicherungsbedingungen:} Es gelten die
Musterbedingungen MB/KK~2009 (Teil~I) sowie die Tarifbedingungen des
Tarifs VitalPlus (Teil~II) in der jeweils gültigen Fassung; eine
Kurzübersicht des Leistungsversprechens enthält Anlage~\ref{sec:anlage1}.

\subsection{Beobachtungseinheiten}\label{sec:beobachtungseinheiten}

Der Tarif umfasst die Beobachtungseinheiten
\begin{center}\small
\begin{tabular}{@{}lll@{}}
\toprule
Beobachtungseinheit & Alter & Kalkulation \\
\midrule
Kinder & 0--15 & natürliche Prämie, keine Alterungsrückstellung \\
männliche Jugendliche & 16--20 & natürliche Prämie, keine Alterungsrückstellung \\
weibliche Jugendliche & 16--20 & natürliche Prämie, keine Alterungsrückstellung \\
Männer & 21--110 & Anwartschaftsdeckungsverfahren \\
Frauen & 21--110 & Anwartschaftsdeckungsverfahren \\
\bottomrule
\end{tabular}
\end{center}
Die Kinder werden geschlechtsunabhängig kalkuliert; die Jugendlichen
werden geschlechtsabhängig in männliche und weibliche Jugendliche
aufgeteilt. Versicherte werden bis zur Vollendung des 16.~Lebensjahres in
der Beobachtungseinheit der Kinder, bis zur Vollendung des
21.~Lebensjahres in der der Jugendlichen geführt (\S\,10 Abs.\,3 KVAV).

\subsection{Altersbestimmung}\label{sec:altersbestimmung}

Das versicherungstechnische Alter ist die Differenz aus laufendem
Kalenderjahr und Geburtsjahr. Die Umstufung zwischen den
Beobachtungseinheiten erfolgt jeweils zum 1.~Januar des Kalenderjahres,
in dem das maßgebliche Alter erreicht wird. Beim Wechsel in die
Beobachtungseinheit der Männer bzw. Frauen (versicherungstechnisches
Alter 21) wird ohne erneute Risikoprüfung der Tarifbeitrag zum
Eintrittsalter 21 gemäß Abschnitt~\ref{sec:beitragsraten} erhoben;
ab diesem Zeitpunkt beginnen Zillmerung und Aufbau der
Alterungsrückstellung.

\subsection{Rahmen-Berechnungsgrundlagen}\label{sec:rahmentb}

Für die folgenden Sachverhalte gelten ergänzend die genannten
Rahmen-Berechnungsgrundlagen (Rahmen-TB) des Unternehmens in den
angegebenen Versionen (fiktive Dokumente):

\begin{center}\small
\begin{tabular}{@{}llll@{}}
\toprule
Rahmen-TB & Version & Inhalt & gilt für \\
\midrule
RTB-AUZ & 2025-1 & Herleitung des aktuariellen Unternehmenszinses & alle Beobachtungseinheiten \\
RTB-Kosten & 2024-2 & Herleitung der Kostensätze aus der Kostenrechnung & alle Beobachtungseinheiten \\
RTB-AWV & 2.1 & Anwartschaftsversicherung (groß/klein) & alle Beobachtungseinheiten \\
\bottomrule
\end{tabular}
\end{center}
